% ==============================================================================
% Chapter 2: Background Knowledge and Related Work
% ==============================================================================
\chapter{Background Knowledge and Related Work}

\section{Background Knowledge}
Several areas of knowledge were necessary for this internship:

\begin{itemize}[label=-]
    \item \textbf{Unity and C\#}: Unity's component-based architecture, ScriptableObjects for data-driven design, and the Animator and physics systems, used to build the game prototype.
    \item \textbf{Object-oriented design and design patterns}: Finite state machines, the Strategy pattern, factories, and object pooling, used to structure enemy behavior and game systems.
    \item \textbf{Software design documentation}: Game Design Documents (GDD), Technical Design Documents (TDD), and UML class diagrams, used to plan the game before implementation.
    \item \textbf{Agile and Scrum}: Sprint-based development with daily meetings, as practiced by the studio.
    \item \textbf{Git workflow and code review}: Branching and pull requests, with every change reviewed by team members before merging.
\end{itemize}

\newpage
\section{Related Work}
\begin{itemize}[label=-]
    \item \textbf{Modern code review}: Bacchelli and Bird \cite{bacchelli2013expectations} studied code review at a large software company and found that review serves purposes beyond finding defects, such as knowledge transfer, team awareness, and finding better solutions. This is relevant to my experience, where review comments helped me learn how to make my code cleaner and more scalable.
    \item \textbf{Problems in game development}: Washburn et al. \cite{washburn2016postmortems} analyzed 155 game development postmortems and identified recurring problems, including scope and schedule issues. This relates to my internship, where project scope had to be defined carefully because the time available was limited.
    \item \textbf{Design patterns in games}: Nystrom \cite{nystrom2014game} describes how patterns such as State, Object Pool, and others apply to game code, including the trade-offs of using them. This relates to the architecture decisions in my project and to my lead's focus on structure.
    \item \textbf{Previous solo project (Mashima)}: My earlier cooperative education project in Japan was developed solo in about two months. Progress was discussed in weekly meetings with my supervisors, but the feedback was positive and did not include technical recommendations, and my code did not go through a formal review process. It serves as a baseline for comparing the professional workflow in this report.
\end{itemize}